\makeabstract
    {
    Investigating Inversion Symmetry Breaking in Few-Layer Nb\textsubscript{3}Cl\textsubscript{8}
    }
    {
    Iris Xiong\textsuperscript{1,2}, Taek Sun Jung\textsuperscript{1}, Liang Wu\textsuperscript{1}
    }
    {
    \textsuperscript{1} Department of Physics \& Astronomy, University of Pennsylvania, Philadelphia, PA, USA\\
    \textsuperscript{2} Department of Physics \& Astronomy, Amherst College, Amherst, MA, USA
    }
    {
    Nb\textsubscript{3}Cl\textsubscript{8} is a breathing-mode kagome 2D van der Waals material. The breathing distortion in its kagome lattice produces alternating large and small Nb\textsubscript{3} trimers that break inversion symmetry within each layer. In multilayer systems, it is theorized to have layer-parity-dependent inversion symmetry, with even layers being centrosymmetric and odd layers noncentrosymmetric. This parity rule is behind many of the material's quantum properties, including topological band structure and out-of-plane polarization. Second harmonic generation (SHG) is a nonlinear optical response that can detect inversion symmetry breaking. We investigated the inversion symmetry of few-layer Nb\textsubscript{3}Cl\textsubscript{8} using polarization-dependent SHG (PSHG) to further understand its layer-parity dependent symmetry.
    }
    {
    Few-layer Nb\textsubscript{3}Cl\textsubscript{8} samples were fabricated by heat assisted mechanical exfoliation. A small bulk flake was thinned using Blue Nitto tape and Scotch Magic tape, then transferred onto a SiO\textsubscript{2}/Si substrate at $100^\circ$C for 3 minutes and cooled to room temperature. Samples were imaged using optical microscopy and atomic force microscopy (AFM). AFM measured the thickness and number of layers $N$ of each Nb\textsubscript{3}Cl\textsubscript{8} exfoliated flake. PSHG probed the inversion symmetry of samples with a 800 nm pulsed laser and a power of 2 mW at room temperature.
    }
    {
    AFM measurements showed samples with a minimum of 2 layers ($N \geq 2$) and thicknesses ranging from $1.2 \pm 0.1$nm ($N = 2$, not shown) to $8.3 \pm 0.2$nm ($N = 12$). Cross and parallel PSHG measurements in sample flake $N = 5$ showed six-fold symmetry, with further analysis using a Fourier fit ($n \leq 6$). This suggests a dominating C3 rotational symmetry that is revealed very likely due to the expected broken inversion symmetry in odd-layered flakes. However, inherent surface effects and possible crystal structure change due to laser damage evidenced by AFM imaging and an observable increased SHG intensity may be alternate sources of broken inversion symmetry.
    }
    {
    Few-layer Nb\textsubscript{3}Cl\textsubscript{8} sample flakes were fabricated through heat-assisted mechanical exfoliation. AFM imaging confirmed the thickness and layer number ($N$), with all measured samples having layers $N \geq 2$. The dominating six-fold symmetry revealed by polarization-dependent second harmonic generation (PSHG) measurements is likely due to broken inversion symmetry in odd layer flakes. In future work, laser damage at the site of SHG should be evaluated with compositional analysis techniques such as Scanning Electron Microscopy with Energy Dispersive X-ray Spectroscopy (SEM-EDS) measurements. Nb\textsubscript{3}Cl\textsubscript{8} samples will also be exfoliated to the monolayer and be investigated with polarization-dependent SHG and Raman spectroscopy measurements.
    }

    \newpage
    